\section{Complexity of hierarchical elimination}\label{sec:complexity}

\subsection{Local elimination fronts}\label{sec:local}

Fix an elimination order inside each $I_j$, with $\partial I_j$ ordered after
the interior. For an interior pivot $v$, let $d_v$ be its number of later
neighbors, interior or retained, in the filled local augmented graph. Define
\begin{equation}\label{eq:phi}
\Phi_1=\sum_j\sum_{v\in I_j}(1+d_v),\qquad
\Phi_2=\sum_j\sum_{v\in I_j}(1+d_v)^2,\qquad
\mathcal B_2=\sum_jb_j^2.
\end{equation}
These quantities depend on the symbolic ordering, not on the numerical values
at a particular iterate. If relaxed supernodes or pattern unions add structural
entries, the padded pattern must be used.

\begin{lemma}[partial factorization cost]\label{lem:local}
All local factors and their retained-front updates can be formed with
\begin{equation}\label{eq:local-cost}
W_{\mathrm{local}}=O(\Phi_2+\mathcal B_2),\qquad M_{\mathrm{local}}=O(\Phi_1+\mathcal B_2).
\end{equation}
With these factors, local condensation and interior recovery cost $O(\Phi_1)$
per right-hand side.
\end{lemma}

\begin{proof}
A scalar pivot with $d_v$ later neighbors normalizes its column and applies a
symmetric rank-one update to its later neighbors, at cost $O((1+d_v)^2)$, and
stores $O(1+d_v)$ factor entries. Summing over interior pivots gives $\Phi_2$ and
$\Phi_1$. Initializing, extracting and storing each retained front adds
$O(b_j^2)$ work and memory. Forward elimination and backward recovery touch each
stored factor entry a constant number of times per right-hand side.
\end{proof}

Let $w_j$ be the maximum number of later \emph{interior} neighbors when
factoring $A_j$. Retained vertices are never local pivots and cannot create
edges between interior vertices, so
\begin{equation}\label{eq:dv}
d_v\le w_j+b_j\quad(v\in I_j),\qquad
\Phi_2\le\sum_jn_j(1+w_j+b_j)^2,\qquad
\Phi_1\le\sum_jn_j(1+w_j+b_j).
\end{equation}
The accumulated measure $\Phi_2$ can be much smaller than a bound that charges
the largest front to every pivot: a few large fronts need not determine the cost
of every variable inside a stage.

\subsection{The reduced chain with a global border}

Partition the reduced system as
\begin{equation}\label{eq:reduced-block}
S=\begin{bmatrix}T&Q\\Q^\top&C_B\end{bmatrix},
\end{equation}
where $T$ is block tridiagonal with $K$ blocks of size at most $p$ and $C_B$ has
order $g$; both $S$ and $T$ are SPD. We use block cyclic
reduction~\cite{hockney1965fast,heller1976some} with a shrinking active set:
each level eliminates alternating active blocks, and the solve visits the levels
in reverse.

\begin{lemma}[reduced solve]\label{lem:reduced}
The system~\eqref{eq:reduced-block} can be factored and prepared with
\begin{equation}\label{eq:reduced-cost}
W_R=O(Kp^3+Kp^2g+Kpg^2+g^3),\qquad M_R=O(Kp^2+Kpg+g^2),
\end{equation}
and every further right-hand side costs $O(Kp^2+Kpg+g^2)$.
\end{lemma}

\begin{proof}
An active block has at most two active neighbors, so its elimination costs
$O(p^3)$ for the dense factorization, triangular solves and neighbor updates.
The active set roughly halves at each level; there are $K-1$ eliminated blocks
and one root, giving $O(Kp^3)$ work, $O(Kp^2)$ storage and $O(Kp^2)$ per solve.
Next, $Z=T^{-1}Q$ needs $g$ chain solves, at cost $O(Kp^2g)$ and storage $O(Kpg)$.
Forming $C_B-Q^\top Z$ costs $O(Kpg^2)$, and factoring it costs $O(g^3)$; the
accumulation can be streamed, so no per-block $g\times g$ update needs to be
stored. A new right-hand side requires one chain solve, a border solve and a
recovery update with $Z$.
\end{proof}

Linear work depends on the shrinking active set. A parallel variant that updates
every block at every level costs $O(Kp^3\log K)$ instead.

\subsection{Main complexity theorem}

\begin{theorem}[hierarchical numerical elimination]\label{thm:main}
Suppose $H_\delta\succ0$, the retained partition satisfies
\cref{def:certificate}, and the local orders are fixed. An exact hierarchical
solver attains
\begin{align}
W_{\mathrm{factor}}&=O\bigl(\Phi_2+\mathcal B_2+Kp^3+Kp^2g+Kpg^2+g^3\bigr),\label{eq:wfactor}\\
M_{\mathrm{numeric}}&=O\bigl(\Phi_1+\mathcal B_2+Kp^2+Kpg+g^2\bigr),\label{eq:mnum}\\
W_{\mathrm{rhs}}&=O\bigl(q_{\mathrm{rhs}}\,(\Phi_1+Kp^2+Kpg+g^2)\bigr).\label{eq:wrhs}
\end{align}
If the Sherman--Morrison preparation~\eqref{eq:sm} is not among the right-hand
sides, add one base solve and $O(m)$ vector work. The case $K=0$ is a dense
solve on $B$.
\end{theorem}

\begin{proof}
\Cref{lem:local} bounds all interior partial factorizations. Accumulating their
fronts into the reduced pattern costs $O(\mathcal B_2)$. \Cref{thm:chain} gives
the structure and positive definiteness needed by \cref{lem:reduced}. Adding the
local and reduced costs gives~\eqref{eq:wfactor} and~\eqref{eq:mnum}, and summing
forward condensation, reduced solves and recovery over the right-hand sides
gives~\eqref{eq:wrhs}.
\end{proof}

\begin{corollary}[linear complexity in the horizon]\label{cor:linear}
If in addition $\sum_jn_j=O(N)$, $K=O(N)$, and $w_j,b_j,p,g$ are bounded
uniformly in $N$, then for fixed $q_{\mathrm{rhs}}$ the numerical factorization,
its storage and the subsequent solves cost $O(N)$.
\end{corollary}

\begin{proof}
By~\eqref{eq:dv}, $\Phi_1,\Phi_2=O(N)$. Every component has at least one vertex,
so $J=O(N)$ and $\mathcal B_2=O(N)$. All remaining terms are linear in $K$.
\end{proof}

The certificate restricts which groups a component touches; it does not bound
the component size. A fixed coarsening span $h$ and a bounded number of rows per
stage supply that bound, provided each component is checked to lie within a
bounded number of stages. Substituting the construction of \cref{thm:spaced}
gives a bound purely in terms of the model parameters.

\begin{corollary}[explicit SOS parameter bound]\label{cor:explicit}
Under \cref{thm:spaced},
\begin{align}
W_{\mathrm{factor}}&=O\Bigl(Nq(hq+2P+g)^2+\tfrac Nh\,(P^3+P^2g+Pg^2)+g^3\Bigr),\label{eq:explicit-w}\\
M_{\mathrm{numeric}}&=O\Bigl(Nq\bigl[hq+2P+g+(2P+g)^2\bigr]+\tfrac Nh\,(P^2+Pg)+g^2\Bigr).\label{eq:explicit-m}
\end{align}
\end{corollary}

\begin{proof}
There are at most $Nq$ interior vertices and $w_j\le n_j-1\le hq-1$. Apply
\eqref{eq:dv}, $b_j\le2P+g$ and $J\le Nq$ to the local terms of \cref{thm:main},
and $K\le N/h$, $p\le P$ to the reduced terms.
\end{proof}

With $q\le\binom{D+2r}{2r}$, \cref{cor:explicit} depends only on the horizon,
the stage and shared variable counts, the relaxation order, the slab span and
the border. It is deliberately conservative, and the front-dependent
\cref{thm:main} is the better predictor in practice. Both make clear that linear
scaling in $N$ comes with a steep dependence on~$r$.

\subsection{Finer accounting}

\paragraph{Unequal interfaces.}
Let $\mathcal P$ be the reduced elimination events, including the root. For
event $e$, let $p_e$ be its pivot-block size and $\beta_e$ the total size of its
(at most two) later neighboring chain blocks. Counting block operations gives
\begin{equation}\label{eq:ragged}
W_T=O\Bigl(\sum_{e\in\mathcal P}\bigl[p_e^3+p_e^2\beta_e+p_e\beta_e^2\bigr]\Bigr),\qquad
M_T=O\Bigl(\sum_{e\in\mathcal P}\bigl[p_e^2+p_e\beta_e\bigr]\Bigr),
\end{equation}
with $O(M_T)$ work per solve. Since $\beta_e\le2p$ and $|\mathcal P|=K$,
\eqref{eq:ragged} implies the chain terms of \cref{lem:reduced} and quantifies
the overhead of padding every interface to the largest width.

\paragraph{The executed supernodal plan.}
For a direct match with the GPU backend, let supernode $s$ have pivot width
$\omega_s$ and panel height $h_s$, and let $\mathcal E_{\mathrm{upd}}$ be the
scheduled source-to-target panel updates. The retained front is an update
target but never a pivot. One numerical instance of a padded plan costs
\begin{equation}\label{eq:panel}
W_{\mathrm{panel}}=O\Bigl[\sum_{s\in\mathrm{pivot}}\bigl(\omega_s^3+(h_s-\omega_s)\omega_s^2\bigr)
+\sum_{(s,t)\in\mathcal E_{\mathrm{upd}}}\omega_s\omega_th_t\Bigr],
\end{equation}
corresponding to POTRF, TRSM and panel updates, with panel storage
$O(\sum_sh_s\omega_s)$. If one symbolic template serves several components,
the numerical cost is multiplied by the number of instances: reusing a template
saves symbolic work but not factorization work.

\subsection{Parallel depth}

Let $D_{\mathrm{loc}}$ be the maximum dependency depth of the local schedules.

\begin{proposition}[chain elimination depth]\label{prop:depth}
Cyclic reduction has $O(\log(K+1))$ forward and backward levels. Executing each
block operation sequentially gives factorization depth $O(p^3\log(K+1))$ and
solve depth $O(p^2\log(K+1))$.
\end{proposition}

\begin{proof}
Each level halves the active set approximately. Blocks eliminated in one level
are nonadjacent and compute their updates independently; each surviving block
receives at most two contributions, combined in a fixed order. Recovery
reverses the levels.
\end{proof}

\begin{corollary}[linear work and logarithmic depth]\label{cor:depth}
Suppose $p,g$, the component sizes and the local arithmetic costs are bounded
uniformly, and $J,K=O(N)$. If independent components run in parallel and
contributions to a common retained entry are combined by balanced reductions,
there is an elimination schedule with
\begin{equation}\label{eq:depth}
W=O(N),\qquad M=O(N),\qquad D=O(\log(N+1)).
\end{equation}
\end{corollary}

\begin{proof}
Local components have constant depth. Combining their contributions has depth
$O(\log(J+1))$, the chain has depth $O(\log(K+1))$ by \cref{prop:depth}, and the
border needs a bounded number of chain solves and reductions. \Cref{thm:main}
supplies the work and storage bounds.
\end{proof}

With greedy scheduling on $P_{\mathrm{proc}}$ processors, the ideal time is
$O(W/P_{\mathrm{proc}}+D)$; bandwidth, synchronization, kernel launches and
occupancy are not modeled. Bounded front width alone does not imply logarithmic
depth, since a component can contain a serial chain of $\Theta(N)$ pivots.

\subsection{Cost of one interior-point iteration}\label{sec:iter-cost}

For PSD block $c$, let $a_{ci}$ be the number of stored lower-triangular
coordinates of $A_i^{(c)}$ and $z_c=\sum_ia_{ci}$.

\begin{proposition}[selector assembly]\label{prop:selector}
Given $W_c$ and precomputed scatter maps, evaluating all incident row pairs of
block $c$ by coordinate pairing requires
\begin{equation}\label{eq:selector}
P_c=\sum_{i\le j}a_{ci}a_{cj}=\tfrac12\Bigl(z_c^2+\sum_ia_{ci}^2\Bigr)
\end{equation}
coordinate-pair evaluations, hence $\Theta(z_c^2)$ operations for a nonempty block.
\end{proposition}

\begin{proof}
Expanding the trace in~\eqref{eq:schur} over two coefficient coordinates uses a
constant number of entries of $W_c$. Expanding $z_c^2$ counts off-diagonal row
pairs twice and diagonal pairs once, which gives the identity, and
$\sum_ia_{ci}^2\le z_c^2$ gives the bounds.
\end{proof}

For a monomial Gram map, each symmetric coordinate belongs to one row, so
$z_c\le k_c(k_c+1)/2$; a localizing polynomial with $\sigma_c$ terms gives
$z_c\le\sigma_ck_c(k_c+1)/2$ and work $O(\sigma_c^2k_c^4)$. The generic dense
path instead costs $O(q_ck_c^3)$ in matrix products per block, so neither path
dominates for all patterns, and the implementation selects per block. With
$t_a$ the number of rows touching free variable $a$, one iteration costs
\begin{equation}\label{eq:witer}
W_{\mathrm{iter}}=O\Bigl(\sum_ck_c^3+\sum_cz_c^2+\sum_at_a^2+\nnz(\mathcal A,E)+\nnz(H_\delta)
+W_{\mathrm{factor}}+W_{\mathrm{rhs}}+m+n_f\Bigr),
\end{equation}
assuming a bounded number of refinement steps, rejected directions and
refactorizations per accepted iteration. Under the conditions of
\cref{cor:linear} and bounded resources per stage, this is $O(N)$.

\paragraph{Outer iterations.}
The total work is $W_{\mathrm{setup}}+\sum_{k=1}^{L}W_{\mathrm{iter},k}$ over
$L$ outer iterations. A short-step path-following method needs
$O(\sqrt\nu\log(1/\epsilon))$ iterations for barrier parameter $\nu$; with
$\nu=O(N)$ this gives $O(N^{3/2}\log(1/\epsilon))$ work beyond
setup~\cite{zhang2020clique}. Our implementation uses a homogeneous self-dual
Mehrotra-type method with safeguards, and we do not claim that it satisfies the
assumptions of that theory. In our experiments the iteration count grows slowly
(23, 24 and 24--25 iterations for $N=10,30,100$ on PushT; \cref{sec:experiments}).

\subsection{Relation to general sparse Cholesky}

\begin{proposition}[representation by an ordering]\label{prop:ordering}
In exact arithmetic, the interior eliminations, the reduced-chain eliminations
and the final border factorization correspond to a permutation under which
ordinary sparse Cholesky performs the same elimination.
\end{proposition}

\begin{proof}
List the interior pivots first, respecting each component's order; components
share no interior edges, so any interleaving is admissible. Then list the chain
pivots by elimination level, expanding blocks into scalar pivots, and finally the
border. By \cref{prop:assoc}, block and scalar elimination in this order produce
the same Schur complements, and every pivot exists because the matrix is SPD.
\end{proof}

A general supernodal or multifrontal solver given this ordering and a
comparable blocking can therefore match the arithmetic bound, and we do not
claim that general sparse Cholesky is asymptotically slower on this matrix
family. The practical question is how to obtain a good ordering and execution
plan from the SOS model. The hierarchy answers it in four ways. Temporal
coefficient supports propose interfaces, the certificate controls fill after
local elimination, pruning and coarsening shrink the reduced system, and
repeated component patterns enable batching. \Cref{sec:cudss} measures the
resulting gain against a general GPU factorization of the same operator.
